\documentclass{beamer}
\usepackage{propositions}

%% Beamer overlay specifications on prop and its items.
%%
%% Three places take one, as in beamer's own environments: an item, an
%% environment's items by default, and the environment itself.  What makes
%% this work is closing the actionenv beamer opens around each item before
%% the next item's own work rather than in the middle of it; otherwise every
%% label in a list comes out as the first item's, the group's end having
%% rolled the variable holding it back.  And the package's counters are
%% registered with \resetcounteronoverlays, or a list numbered (1), (2) on
%% its first overlay comes out (3), (4) on its second.

\begin{document}

\begin{frame}{One item at a time}
Two overlays, numbered (1) and (2) on both of them.
\begin{prop}
  \item<1-> Alpha
  \item<2-> Beta
\end{prop}
\end{frame}

\begin{frame}{A specification on either side of the optional argument}
Both orders are beamer's; both must work.
\begin{prop}
  \item<1->[Premise] Alpha
  \item[Conclusion]<2-> Beta
\end{prop}
\end{frame}

\begin{frame}{The environment's default, and keys after it}
\begin{prop}[<+->]
  \item Alpha
  \item Beta
\end{prop}
\begin{prop}[<+->, align=runin]
  \item Gamma
  \item Delta
\end{prop}
\end{frame}

\begin{frame}{A specification on the environment itself}
Nothing on the first overlay; the whole list on the second.
\begin{prop}<2->
  \item Alpha
  \item Beta
\end{prop}
\end{frame}

\begin{frame}{The two composing}
The environment is \texttt{<2-4>}, the items \texttt{<1->} and \texttt{<3->}.
An item appears where both are satisfied: nothing, then Alpha, then both,
then both.
\begin{prop}<2-4>
  \item<1-> Alpha
  \item<3-> Beta
\end{prop}
\end{frame}

\begin{frame}{Nesting, and every alignment}
\begin{prop}
  \item<1-> Outer
  \begin{prop}
    \item<2-> Inner
  \end{prop}
  \item<3-> Outer again
\end{prop}
\begin{prop}[<+->, align=nextline]\item Nextline one \item Nextline two\end{prop}
\begin{prop}[<+->, align=flush]\item Flush one \item Flush two\end{prop}
\begin{prop}[<+->, align=center]\item Centre one \item Centre two\end{prop}
\end{frame}

\begin{frame}{inlineprop, the same three}
An \texttt{inlineprop} calls no \texttt{\string\item} of beamer's, so it opens
and closes the overlay group itself.  It is the same \texttt{actionenv} a
\texttt{prop} item gets: an item not yet reached holds its width, and a label
inside one survives.
Consider: \begin{inlineprop}\item<1-> one,\label{b:i} \item<2-> two,
\item<3-> three.\label{b:iii}\end{inlineprop} And after.
\end{frame}

\begin{frame}{inlineprop, environment default and environment specification}
\begin{inlineprop}[<+->]\item one, \item two.\end{inlineprop}

Before. \begin{inlineprop}<2->\item one, \item two.\end{inlineprop} After.
\end{frame}

\begin{frame}{inlineprop nested, in a prop item and in itself}
\begin{prop}
  \item<1-> Outer, with \begin{inlineprop}\item<2-> in one,
  \item<3-> in two.\end{inlineprop} done.
\end{prop}
\begin{inlineprop}\item<1-> outer, \begin{inlineprop}\item<2-> inner one,
\item<3-> inner two,\end{inlineprop} \item<4-> outer two.\end{inlineprop}
\end{frame}

\begin{frame}{Cross-references and \texttt{\string\pause}}
\begin{prop}
  \item Alpha \label{b:a} \pause
  \item Beta \label{b:b}
\end{prop}
References: \ref{b:a}, \ref{b:b}, \ref{b:i}, \ref{b:iii}.
\end{frame}

\end{document}
